\chapter{2009年大纲I卷（文）}

\section{单选题}

\begin{problem}
\(\sin 585^{\circ}\) 的值为（\quad）

\choices
    {\(-\frac{\sqrt{2}}{2}\)}
    {\( \frac{\sqrt{2}}{2} \)}
    {\(-\frac{\sqrt{3}}{2}\)}
    {\(\frac{\sqrt{3}}{2}\)}
\end{problem}

\begin{problem}
设集合 \(A = \{4, 5, 7, 9\}\)，\(B = \{3, 4, 7, 8, 9\}\)，全集 \(U = A \cup B\)，则集合 \(\complement_U (A \cap B)\) 中的元素共有（\quad）

\choices
    {3 个}
    {4 个}
    {5 个}
    {6 个}
\end{problem}

\begin{problem}
不等式 \(\left|\frac{x + 1}{x - 1}\right| < 1\) 的解集为（\quad）

\choices
    {\(\{x\mid 0 < x < 1\} \cup \{x\mid x > 1\}\)}
    {\( \{x \mid 0 < x < 1\} \)}
    {\( \{x \mid -1 < x < 0\} \)}
    {\( \{x \mid x < 0\} \)}
\end{problem}

\begin{problem}
已知 \(\tan \alpha = 4, \cot \beta = \frac{1}{3}\)，则 \(\tan (\alpha + \beta) =\)（\quad）

\choices
    {\(\frac{7}{11}\)}
    {\(-\frac{7}{11}\)}
    {\(\frac{7}{13}\)}
    {\(-\frac{7}{13}\)}
\end{problem}

\begin{problem}
设双曲线 \(\frac{x^2}{a^2} - \frac{y^2}{b^2} = 1 (a > 0, b > 0)\) 的渐近线与抛物线 \(y = x^2 + 1\) 相切. 则该双曲线的离心率等于（\quad）

\choices
    {\(\sqrt{3}\)}
    {2}
    {\(\sqrt{5}\)}
    {\(\sqrt{6}\)}
\end{problem}

\begin{problem}
已知函数 \( f(x) \) 的反函数为 \( g(x) = 1 + 2\lg x \)（\(x > 0\)），则 \( f(1) + g(1) = \)（\quad）

\choices
    {0}
    {1}
    {2}
    {4}
\end{problem}

\begin{problem}
甲组有 5 名男同学，3 名女同学；乙组有 6 名男同学，2 名女同学. 若从甲，乙两组中各选出 2 名同学，则选出的 4 人中恰有 1 名女同学的不同选法共有（\quad）

\choices
    {150 种}
    {180 种}
    {300 种}
    {345 种}
\end{problem}

\begin{problem}
设非零向量 \(a, b, c\) 满足 \(|a| = |b| = |c|, a + b = c\)，则 \(\langle a, b \rangle =\)（\quad）

\choices
    {\(150^{\circ}\)}
    {\(120^{\circ}\)}
    {\(60^{\circ}\)}
    {\(30^{\circ}\)}
\end{problem}

\begin{problem}
已知三棱柱 \(ABC - A_{1}B_{1}C_{1}\) 的侧棱与底面边长都相等，\(A_{1}\) 在底面 \(ABC\) 上的射影为 \(BC\) 的中点，则异面直线 \(AB\) 与 \(CC_{1}\) 所成的角的余弦值为（\quad）

\choices
    {\(\frac{\sqrt{3}}{4}\)}
    {\(\frac{\sqrt{5}}{4}\)}
    {\(\frac{\sqrt{7}}{4}\)}
    {\(\frac{3}{4}\)}
\end{problem}

\begin{problem}
如果函数 \( y = 3\cos (2x + \varphi) \) 的图象关于点 \( \left(\frac{4\pi}{3},0\right) \) 中心对称，那么 \( |\varphi| \) 的最小值为（\quad）

\choices
    {\(\frac{\pi}{6}\)}
    {\( \frac{\pi}{4} \)}
    {\( \frac{\pi}{3} \)}
    {\( \frac{\pi}{2} \)}
\end{problem}

\begin{problem}
已知二面角 \(\alpha - l - \beta\) 为 \(60^{\circ}\)，动点 \(P, Q\) 分别在面 \(\alpha, \beta\) 内，\(P\) 到 \(\beta\) 的距离为 \(\sqrt{3}, Q\) 到 \(\alpha\) 的距离为 \(2\sqrt{3}\)，则 \(P, Q\) 两点之间距离的最小值为（\quad）

\choices
    {\(\sqrt{2}\)}
    {2}
    {\(2\sqrt{3}\)}
    {4}
\end{problem}

\begin{problem}
已知椭圆 \(C: \frac{x^2}{2} + y^2 = 1\) 的右焦点为 \(F\)，右准线 \(l\)，点 \(A \in l\)，线段 \(AF\) 交 \(C\) 于点 \(B\). 若 \(\overrightarrow{FA} = 3\overrightarrow{FB}\)，则 \(|\overrightarrow{AF}| =\)（\quad）

\choices
    {\(\sqrt{2}\)}
    {2}
    {\(\sqrt{3}\)}
    {3}
\end{problem}

\begin{problem}
已知椭圆 \(C: \frac{x^2}{2} + y^2 = 1\) 的右焦点为 \(F\)，右准线为 \(l\)，点 \(A \in l\)，线段 \(AF\) 交 \(C\) 于点 \(B\). 若 \(\overrightarrow{FA} = 3\overrightarrow{FB}\)，则 \(|AF| =\)（\quad）

\choices
    {\(\sqrt{2}\)}
    {2}
    {\(\sqrt{3}\)}
    {3}
\end{problem}

\section{填空题}

\begin{problem}
\((x - y)^{10}\) 的展开式中，\(x^{7}y^{3}\) 的系数与 \(x^{3}y^{7}\) 的系数之和等于\fillinblank{}.
\end{problem}

\begin{problem}
设等差数列 \(\{a_{n}\}\) 的前 \(n\) 项和为 \(S_{n}\)，若 \(S_{9} = 72\)，则 \(a_{2} + a_{4} + a_{9} = \)\fillinblank{}.
\end{problem}

\begin{problem}
已知 \(OA\) 为球 \(O\) 的半径，过 \(OA\) 的中点 \(M\) 且垂直于 \(OA\) 的平面截球面得到圆 \(M\). 若圆 \(M\) 的面积为 \(3\pi\)，则球 \(O\) 的表面积等于\fillinblank{}.
\end{problem}

\begin{problem}
若直线 \(m\) 被两平行线 \(l_{1}: x - y + 1 = 0\) 与 \(l_{2}: x - y + 3 = 0\) 所截得的线段的长为 \(2\sqrt{2}\)，则 \(m\) 的倾斜角可以是 \circled{1} \(15^{\circ}\) \circled{2} \(30^{\circ}\) \circled{3} \(45^{\circ}\) \circled{4} \(60^{\circ}\) \circled{5} \(75^{\circ}\)，其中正确答案的序号是\fillinblank{}. (写出所有正确答案的序号)
\end{problem}

\section{解答题}

\begin{problem}
设等差数列 \(\{a_{n}\}\) 的前 \(n\) 项和为 \(S_{n}\)，公比是正数的等比数列 \(\{b_{n}\}\) 的前 \(n\) 项和为 \(T_{n}\)，已知 \(a_{1} = 1, b_{1} = 3, a_{3} + b_{3} = 17, T_{3} - S_{3} = 12\)，求 \(\{a_{n}\}, \{b_{n}\}\) 的通项公式.
\end{problem}

\begin{problem}
在 \(\triangle ABC\) 中，内角 \(A\)，\(B\)，\(C\) 的对边长分别为 \(a\)，\(b\)，\(c\)，已知 \(a^{2} - c^{2} = 2b\)，且 \(\sin B = 4\cos A \sin C\)，求 \(b\).
\end{problem}

\begin{problem}
如图，四棱锥 \(S - ABCD\) 中，底面 \(ABCD\) 为矩形，\(SD \perp\) 底面 \(ABCD\)，\(AD = \sqrt{2}, DC = SD = 2\)，点 \(M\) 在侧棱 \(SC\) 上，\(\angle ABM = 60^{\circ}\).

\begin{enumerate}
\item 证明：\(M\) 是侧棱 \(SC\) 的中点；
\item 求二面角 \(S - AM - B\) 的大小.
\end{enumerate}

\bitmapfigure[max width=0.42\linewidth,max totalheight=4.8cm]{img/2009/syllabus_paper_1_liberal/q20_fig1.png}

\end{problem}

\begin{problem}
甲，乙二人进行一次围棋比赛，约定先胜3局者获得这次比赛的胜利. 比赛结束，假设在一局中，甲获胜的概率为0.6，乙获胜的概率为0.4，各局比赛结果相互独立，已知前 2 局中，甲，乙各胜 1 局.

\begin{enumerate}
\item 求再赛 2 局结束这次比赛的概率；
\item 求甲获得这次比赛胜利的概率.
\end{enumerate}
\end{problem}

\begin{problem}
已知函数 \( f(x) = x^4 - 3x^2 + 6 \).

\begin{enumerate}
\item 讨论 \( f(x) \) 的单调性；
\item 设点 \(P\) 在曲线 \(y = f(x)\) 上，若该曲线在点 \(P\) 处的切线 \(l\) 通过坐标原点，求 \(l\) 的方程.
\end{enumerate}
\end{problem}

\begin{problem}
如图，已知抛物线 \(E: y^{2} = x\) 与圆 \(M: (x - 4)^{2} + y^{2} = r^{2} (r > 0)\) 相交于 \(A, B, C, D\) 四个点.

\begin{enumerate}
\item 求 \(r\) 的取值范围；
\item 当四边形 \(ABCD\) 的面积最大时，求对角线 \(AC\)，\(BD\) 的交点 \(P\) 坐标.
\end{enumerate}

\bitmapfigure[max width=0.42\linewidth,max totalheight=4.8cm]{img/2009/syllabus_paper_1_liberal/q23_fig1.png}

\end{problem}
